\chapter{Conception et implémentation de la solution}
\label{chap:conception}

Ce chapitre décrit la mise en œuvre concrète du pipeline présenté au
chapitre~\ref{chap:etatart}~: l'environnement technique retenu, un
raffinement apporté pour réduire la latence de détection, les
résultats obtenus canal par canal, leur validation par confirmation
multi-canal, et le tableau de bord 3D permettant de les explorer.

\clearpage
\section{Environnement et outils techniques}

Le pipeline est entièrement implémenté en Python, au sein d'un
notebook Jupyter unique couvrant l'ensemble du traitement, de
l'ingestion des fichiers bruts à l'export des résultats
(tableau~\ref{tab:stack}). Le calcul de la distance de
Damerau-Levenshtein, potentiellement coûteux s'il devait être répété
pour chacune des 10 000 grilles et 62 jours, s'appuie sur la
bibliothèque \texttt{rapidfuzz}, qui implémente cette distance en C++
avec une interface Python, plutôt que sur une implémentation naïve en
Python pur.

\begin{table}[htbp]
  \centering
  \caption{Environnement et outils techniques}
  \label{tab:stack}
  \begin{tabular}{ll}
    \toprule
    \textbf{Usage} & \textbf{Technologie} \\
    \midrule
    Traitement des données        & Python, pandas, numpy \\
    Distance d'édition            & rapidfuzz (Damerau-Levenshtein) \\
    Géométrie / cartes statiques  & geopandas, folium, matplotlib \\
    Analyse exploratoire          & Jupyter Notebook \\
    Dashboard 3D                  & deck.gl, MapLibre GL JS, JavaScript vanilla \\
    \bottomrule
  \end{tabular}
\end{table}

L'encodage d'une série de trafic horaire en niveaux (section
\ref{sec:methodo} du chapitre précédent) est implémenté par une
fonction vectorisée, appliquée indépendamment à chaque grille~: la
médiane de la grille est calculée une fois, puis chaque valeur
horaire est comparée aux seuils de la nomenclature par recherche
dichotomique (\texttt{numpy.searchsorted}), ce qui évite une boucle
explicite sur les 24 heures de chaque jour.

\begin{lstlisting}[language=Python, basicstyle=\small\ttfamily, frame=single, caption={Encodage vectorisé d'une série horaire en niveaux A-F}, label={lst:encode}]
LEVEL_THRESHOLDS = [0.43, 0.63, 1.00, 1.35, 2.10]
ALPHABET = "ABCDEF"

def encode_levels(values, thresholds=LEVEL_THRESHOLDS):
    """Encode le trafic horaire d'une grille en niveaux 0..5,
    relatifs a sa mediane."""
    median = values.median()
    if median == 0:
        return pd.Series(0, index=values.index)
    ratio = values / median
    return np.searchsorted(thresholds, ratio, side="right")
\end{lstlisting}

\section{Réduction de la latence de détection~: fenêtre glissante}

La méthode décrite au chapitre précédent détecte une anomalie une
fois par jour calendaire~: il faut attendre minuit pour disposer de
la signature complète et pouvoir la comparer à sa référence, ce qui
introduit potentiellement jusqu'à 24 heures de latence entre le début
d'un comportement anormal et sa détection. Un raffinement a donc été
apporté~: plutôt que de recalculer la signature une fois par jour, on
la recalcule \emph{à chaque heure}, sur une fenêtre glissante des 24
dernières heures (tableau~\ref{tab:sliding}).

\begin{table}[htbp]
  \centering
  \caption{Paramètres de la détection par fenêtre glissante}
  \label{tab:sliding}
  \begin{tabular}{lll}
    \toprule
    \textbf{Paramètre} & \textbf{Ce que c'est} & \textbf{Valeur retenue} \\
    \midrule
    Longueur de la fenêtre    & Nombre d'heures couvertes par la signature & 24 heures \\
    Fréquence d'évaluation    & Cadence de recalcul de la distance         & 1 fois par heure \\
    \bottomrule
  \end{tabular}
\end{table}

Ce changement ne remet pas en cause la construction du profil de
référence~: celui-ci n'est pas stocké comme une séquence associée à
\emph{un jour}, mais case par case, par \emph{(grille, jour de
semaine, heure)}. Une fenêtre chevauchant par exemple mardi soir et
mercredi matin se compare donc simplement à deux portions de
référence différentes, sans recalcul. Ce passage d'environ 52
distances par grille (une par jour) à plus d'un millier (une par
heure) améliore également l'estimation de l'écart type par grille, et
donc la fiabilité du score z.

Une fenêtre de 24 heures reste toutefois peu sensible à une anomalie
très brève~: un pic isolé d'une seule heure ne modifie qu'une lettre
sur 24 dans la signature, rarement suffisant pour franchir le seuil.
Une seconde fenêtre, plus courte (6 heures), tourne donc en parallèle
avec sa propre référence et son propre seuil, afin de ne pas perdre
les pics brefs~; une alerte confirmée par les deux fenêtres est
considérée comme plus fiable qu'une alerte isolée -- le même principe
que la confirmation multi-canal de DiNATrAX~\parencite{maudoux2024dinatrax}
et de ce travail (section~\ref{sec:multicanal}), appliqué ici à deux
échelles temporelles plutôt qu'à deux canaux. Trois points ont
nécessité une attention particulière à l'implémentation~: le
regroupement des alertes consécutives déclenchées par une même
anomalie prolongée (on ne conserve que le front montant, c'est-à-dire
le premier franchissement du seuil jusqu'au retour en dessous), la
gestion des grilles sans variance historique (écart type nul, pour
lesquelles on bascule sur un seuil en valeur absolue), et le fait que
la référence reste calculée hors ligne~: en production, chaque
évaluation horaire ne coûte qu'une comparaison de deux séquences de
24 caractères, quasi instantanée.

Le tableau~\ref{tab:latence} illustre le gain obtenu sur l'anomalie la
plus marquée détectée sur la période -- la grille 1421 le 1\textsuperscript{er}
janvier 2014 (Nouvel An), avec un score z de bloc journalier de 5,00.

\begin{table}[htbp]
  \centering
  \caption{Gain de latence sur l'anomalie du Nouvel An (grille 1421, 1\textsuperscript{er} janvier 2014)}
  \label{tab:latence}
  \begin{tabular}{lc}
    \toprule
    \textbf{Méthode} & \textbf{Latence de détection} \\
    \midrule
    Bloc journalier (référence, chapitre~\ref{chap:etatart}) & jusqu'à 24 h (détectée à minuit) \\
    Fenêtre glissante (premier franchissement à 12 h) & 13 h \\
    \midrule
    \textbf{Gain} & \textbf{11 h} \\
    \bottomrule
  \end{tabular}
\end{table}

\section{Résultats de la détection par canal}

Le pipeline complet -- encodage, distance de Damerau-Levenshtein,
score z, seuil à $z>3$ -- a été exécuté indépendamment sur les trois
canaux disponibles, sur l'ensemble de la période couverte (62 jours,
10 000 grilles). Le tableau~\ref{tab:resultats-canal} résume, pour
chaque canal, le nombre de journées-grille détectées comme anormales
et leur taux de correspondance avec le calendrier de repère construit
manuellement (section~\ref{sec:calendrier}).

\begin{table}[htbp]
  \centering
  \caption{Résultats de la détection par canal}
  \label{tab:resultats-canal}
  \begin{tabular}{lccc}
    \toprule
    \textbf{Canal} & \textbf{Anomalies détectées} & \textbf{Part des journées-grille} & \textbf{Correspondance calendaire} \\
    \midrule
    Internet & 2 484 & 0,48~\% & 86,8~\% \\
    SMS      & 4 755 & 0,91~\% & 92,8~\% \\
    Appels   & 6 468 & 1,24~\% & 97,6~\% \\
    \bottomrule
  \end{tabular}
\end{table}

Ces résultats appellent deux remarques. D'abord, les canaux SMS et
appels, plus discrets et plus dispersés que le trafic Internet,
produisent davantage d'anomalies en proportion, mais avec un taux de
correspondance calendaire plus élevé~: un écart marqué sur un signal
naturellement stable est plus significatif que le même écart sur un
signal déjà bruité. Ensuite, la calibration des seuils de
nomenclature sur les percentiles réels du trafic (plutôt que sur des
valeurs rondes) réduit le nombre d'anomalies détectées de moitié tout
en améliorant leur taux de correspondance calendaire (86,8~\% contre
83,9~\% pour des seuils ronds sur le canal Internet), confirmant que
cette calibration élimine surtout des faux positifs.

\section{Confirmation multi-canal et validation calendaire}
\label{sec:multicanal}

\subsection{Calendrier de repère}
\label{sec:calendrier}

Un calendrier de repère a été construit manuellement pour la période
étudiée, combinant jours fériés italiens et matchs à domicile joués à
San Siro -- stade partagé par l'AC Milan et l'Inter
(tableau~\ref{tab:calendrier}). Deux matchs (Milan-Ajax le 11
décembre et Milan-Roma le 16 décembre) tombent dans une période sans
données (fichiers manquants du 11 au 20 décembre) et sont conservés
pour la complétude du calendrier, sans pouvoir être confrontés à une
détection.

\begin{table}[htbp]
  \centering
  \caption{Calendrier de repère (jours fériés et matchs à San Siro)}
  \label{tab:calendrier}
  \small
  \begin{tabular}{ll}
    \toprule
    \textbf{Date} & \textbf{Événement} \\
    \midrule
    01/11/2013 & Ognissanti (Toussaint, jour férié) \\
    02/11/2013 & Milan -- Fiorentina (Serie A) \\
    09/11/2013 & Inter -- Livorno (Serie A) \\
    23/11/2013 & Milan -- Genoa (Serie A) \\
    01/12/2013 & Inter -- Sampdoria (Serie A) \\
    04/12/2013 & Inter -- Trapani (Coppa Italia) \\
    07/12/2013 & Sant'Ambrogio (jour férié local, patron de Milan) \\
    08/12/2013 & Immacolata Concezione (jour férié) + Inter -- Parma \\
    11/12/2013 & Milan -- Ajax (Ligue des champions)$^{*}$ \\
    16/12/2013 & Milan -- Roma (Serie A)$^{*}$ \\
    22/12/2013 & Derby Milan -- Inter (Serie A) \\
    24/12/2013 & Veille de Noël \\
    25/12/2013 & Noël (jour férié) \\
    26/12/2013 & Santo Stefano (jour férié) \\
    01/01/2014 & Nouvel An (jour férié) \\
    \bottomrule
  \end{tabular}
  \\[1mm]
  {\footnotesize $^{*}$ Dans la période sans données (11--20 décembre).}
\end{table}

\subsection{Apport de la confirmation multi-canal}

Sur l'ensemble des trois canaux, 11 440 journées-grille distinctes
sont signalées comme anormales par au moins un canal. Seules 2 062
d'entre elles (18~\%) le sont \emph{simultanément} par au moins deux
canaux -- ce sous-ensemble atteint un taux de correspondance
calendaire de 99,2~\% (2 045 sur 2 062), contre 86,8~\% à 97,6~\% pour
une anomalie détectée sur un seul canal (figure~\ref{fig:multicanal}).
Une anomalie confirmée sur plusieurs canaux constitue donc un signal
nettement plus fiable qu'une anomalie isolée.

\begin{figure}[htbp]
  \centering
  \begin{tikzpicture}
    \begin{axis}[
      ybar, bar width=22pt,
      width=0.75\textwidth, height=6cm,
      symbolic x coords={Internet,SMS,Appels,Confirmé ($\geq$2 canaux)},
      xtick=data,
      x tick label style={font=\footnotesize, align=center},
      ymin=0, ymax=105,
      ylabel={correspondance calendaire (\%)},
      nodes near coords,
      nodes near coords style={font=\footnotesize},
      enlarge x limits=0.15,
    ]
      \addplot[fill=parisneutral] coordinates {
        (Internet,86.8) (SMS,92.8) (Appels,97.6) (Confirmé ($\geq$2 canaux),99.2)
      };
    \end{axis}
  \end{tikzpicture}
  \caption{Taux de correspondance calendaire, par canal isolé puis pour les anomalies confirmées sur au moins deux canaux}
  \label{fig:multicanal}
\end{figure}

\section{Tableau de bord 3D interactif}

Les résultats sont explorables via un tableau de bord 3D
(figure~\ref{fig:dashboard}), développé en JavaScript avec deck.gl et
MapLibre GL JS, et chargé à la demande depuis un serveur HTTP local.
La carte représente les 10 000 grilles de Milan sous forme de
colonnes 3D dont la hauteur et la couleur encodent l'intensité du
trafic à l'heure sélectionnée, les grilles en anomalie ce jour-là
étant mises en évidence par une teinte distincte. L'interface propose
un sélecteur de canal (Internet, SMS, appels), une navigation
jour/heure avec lecture automatique, un mode de comparaison au profil
normal (silhouette superposée), ainsi qu'un histogramme du nombre
d'anomalies par jour permettant de naviguer directement vers les
journées les plus atypiques.

\begin{figure}[htbp]
  \centering
  \includegraphics[width=\textwidth]{dashboard-screenshot.jpg}
  \caption{Tableau de bord 3D~: vue du 1\textsuperscript{er} novembre 2013 à 12h, jour de la Toussaint -- 13 grilles en anomalie ce jour-là (en rouge)}
  \label{fig:dashboard}
\end{figure}

\section{Discussion des résultats}

Les résultats obtenus répondent aux objectifs fixés au
chapitre~\ref{chap:contexte}~: la méthode détecte, sans données
étiquetées, un ensemble restreint et interprétable de journées
atypiques par grille, dont une large majorité correspond à un
événement calendaire connu -- et cette correspondance s'améliore
franchement dès lors qu'on exige une confirmation sur plusieurs
canaux. Le raffinement par fenêtre glissante montre par ailleurs
qu'une contrainte opérationnelle (réduire la latence de détection)
peut être satisfaite sans modifier le cœur de la méthode, en jouant
uniquement sur la fréquence de réévaluation. Les anomalies non
expliquées par le calendrier de repère (13,2~\% à 2,4~\% selon le
canal, et moins de 1~\% pour les anomalies confirmées) restent le
principal point ouvert~: sans vérité terrain, il n'est pas possible de
distinguer, parmi elles, un événement réel non répertorié d'un faux
positif résiduel.
